
\documentclass{article}
\usepackage{colortbl}
\usepackage{makecell}
\usepackage{multirow}
\usepackage{supertabular}

\begin{document}

\newcounter{utterance}

\centering \large Interaction Transcript for game `cladder', experiment `full\_v1.5\_default', episode 5930 with Qwen3.5{-}4B.
\vspace{24pt}

{ \footnotesize  \setcounter{utterance}{1}
\setlength{\tabcolsep}{0pt}
\begin{supertabular}{c@{$\;$}|p{.15\linewidth}@{}p{.15\linewidth}p{.15\linewidth}p{.15\linewidth}p{.15\linewidth}p{.15\linewidth}}
    \# & \multicolumn{2}{c}{Player} && \multicolumn{2}{c}{Game Master} \\
    \hline

    \theutterance \stepcounter{utterance}  
    & & & \multicolumn{4}{p{0.6\linewidth}}{
        \cellcolor[rgb]{0.9,0.9,0.9}{
            \makecell[{{p{\linewidth}}}]{
                \texttt{\tiny{[P1$\langle$GM]}}
                \texttt{You are an expert in causal inference. The following question is not a typical commonsense query, but rather a meticulously designed question created by a professor specializing in causal inference, intended to assess the students' mastery of the course content.} \\
\texttt{Imagine a self{-}contained, hypothetical world with only the following conditions, and without any unmentioned factors or causal relationships: Drinking coffee has a direct effect on wife and alarm clock. Wife has a direct effect on alarm clock. We know that drinking coffee causes alarm set by wife. drinking coffee or alarm set by wife causes ringing alarm. Would the alarm rings the next morning if drinking coffee instead of not drinking coffee?} \\
\texttt{Start your answer with 'Yes' or 'No', followed by additional reasoning or evidence to support your explanation.} \\
            }
        }
    }
     \\ \\

    \theutterance \stepcounter{utterance}  
    & & \multicolumn{4}{p{0.6\linewidth}}{
        \cellcolor[rgb]{0.95,0.95,0.95}{
            \makecell[{{p{\linewidth}}}]{
                \texttt{\tiny{[P1$\leftrightarrow$P]}}
                \texttt{{[}DRAFT{]}} \\
\texttt{Yes} \\
\\ 
\texttt{**Reasoning:**} \\
\\ 
\texttt{To determine the outcome, we must analyze the causal structure defined in the hypothetical world:} \\
\\ 
\texttt{1.  **Identify the Intervention:** The question asks for the effect of the intervention \$do(\textbackslash{}text\{Drinking Coffee\})\$. This means we are forcing the state of "Drinking Coffee" to be true, regardless of any other causes.} \\
\texttt{2.  **Trace the Causal Pathways:**} \\
\texttt{    *   **Path A:** \$\textbackslash{}text\{Drinking Coffee\} \textbackslash{}rightarrow \textbackslash{}text\{Wife\} \textbackslash{}rightarrow \textbackslash{}text\{Alarm (set by wife)\} \textbackslash{}rightarrow \textbackslash{}text\{Ringing Alarm\}\$.} \\
\texttt{    *   **Path B:** \$\textbackslash{}text\{Drinking Coffee\} \textbackslash{}rightarrow \textbackslash{}text\{Alarm (set by wife)\} \textbackslash{}rightarrow \textbackslash{}text\{Ringing Alarm\}\$ (Direct effect mentioned).} \\
\texttt{    *   **Path C:** \$\textbackslash{}text\{Drinking Coffee\} \textbackslash{}rightarrow \textbackslash{}text\{Alarm (set by wife)\} \textbackslash{}rightarrow \textbackslash{}text\{Ringing Alarm\}\$ (Via the "drinking coffee or alarm set by wife" condition).} \\
\texttt{3.  **Analyze the "Or" Condition:** The condition states that "drinking coffee or alarm set by wife causes ringing alarm." In causal logic, if either component of an "OR" condition is true, the effect occurs. Since the intervention sets "Drinking Coffee" to true, the condition "Drinking Coffee OR Alarm set by wife" is satisfied.} \\
\texttt{4.  **Evaluate Confounding/Backdoor Paths:** The prompt explicitly states there are "no unmentioned factors or causal relationships." Therefore, we do not need to worry about confounding variables (like the husband's schedule or the time of day) that might block the causal chain or introduce noise. The system is closed.} \\
\texttt{5.  **Conclusion:** Since "Drinking Coffee" is a direct cause of the "Alarm set by wife" (both directly and via the wife), and the "Alarm set by wife" (along with drinking coffee) is a sufficient cause for the "Ringing Alarm," the intervention guarantees the alarm rings.} \\
\\ 
\texttt{Therefore, under the strict conditions provided, the alarm will ring.} \\
            }
        }
    }
    & \\ \\

    \theutterance \stepcounter{utterance}  
    & & \multicolumn{4}{p{0.6\linewidth}}{
        \cellcolor[rgb]{0.95,0.95,0.95}{
            \makecell[{{p{\linewidth}}}]{
                \texttt{\tiny{[P1$\leftrightarrow$P]}}
                \texttt{{[}SELF{-}FEEDBACK{]}} \\
\texttt{FORMAT: ok} \\
\texttt{RULES: ok} \\
\texttt{STRATEGY: ok} \\
\texttt{REASON: The answer correctly identifies the causal chain and provides a logically sound justification for the intervention's outcome.} \\
            }
        }
    }
    & \\ \\

    \theutterance \stepcounter{utterance}  
    & \multicolumn{4}{p{0.6\linewidth}}{
        \cellcolor[rgb]{0.561,0.737,0.561}{
            \makecell[{{p{\linewidth}}}]{
                \texttt{\tiny{[P1$\rangle$GM]}}
                \texttt{Yes} \\
\\ 
\texttt{**Reasoning:**} \\
\\ 
\texttt{To determine the outcome, we must analyze the causal structure defined in the hypothetical world:} \\
\\ 
\texttt{1.  **Identify the Intervention:** The question asks for the effect of the intervention \$do(\textbackslash{}text\{Drinking Coffee\})\$. This means we are forcing the state of "Drinking Coffee" to be true, regardless of any other causes.} \\
\texttt{2.  **Trace the Causal Pathways:**} \\
\texttt{    *   **Path A:** \$\textbackslash{}text\{Drinking Coffee\} \textbackslash{}rightarrow \textbackslash{}text\{Wife\} \textbackslash{}rightarrow \textbackslash{}text\{Alarm (set by wife)\} \textbackslash{}rightarrow \textbackslash{}text\{Ringing Alarm\}\$.} \\
\texttt{    *   **Path B:** \$\textbackslash{}text\{Drinking Coffee\} \textbackslash{}rightarrow \textbackslash{}text\{Alarm (set by wife)\} \textbackslash{}rightarrow \textbackslash{}text\{Ringing Alarm\}\$ (Direct effect mentioned).} \\
\texttt{    *   **Path C:** \$\textbackslash{}text\{Drinking Coffee\} \textbackslash{}rightarrow \textbackslash{}text\{Alarm (set by wife)\} \textbackslash{}rightarrow \textbackslash{}text\{Ringing Alarm\}\$ (Via the "drinking coffee or alarm set by wife" condition).} \\
\texttt{3.  **Analyze the "Or" Condition:** The condition states that "drinking coffee or alarm set by wife causes ringing alarm." In causal logic, if either component of an "OR" condition is true, the effect occurs. Since the intervention sets "Drinking Coffee" to true, the condition "Drinking Coffee OR Alarm set by wife" is satisfied.} \\
\texttt{4.  **Evaluate Confounding/Backdoor Paths:** The prompt explicitly states there are "no unmentioned factors or causal relationships." Therefore, we do not need to worry about confounding variables (like the husband's schedule or the time of day) that might block the causal chain or introduce noise. The system is closed.} \\
\texttt{5.  **Conclusion:** Since "Drinking Coffee" is a direct cause of the "Alarm set by wife" (both directly and via the wife), and the "Alarm set by wife" (along with drinking coffee) is a sufficient cause for the "Ringing Alarm," the intervention guarantees the alarm rings.} \\
\\ 
\texttt{Therefore, under the strict conditions provided, the alarm will ring.} \\
            }
        }
    }
    & & \\ \\

    \theutterance \stepcounter{utterance}  
    & & & \multicolumn{2}{p{0.3\linewidth}}{
        \cellcolor[rgb]{0.9,0.9,0.9}{
            \makecell[{{p{\linewidth}}}]{
                \texttt{\tiny{[GM$|$GM]}}
                \texttt{Yes} \\
            }
        }
    }
    & & \\ \\

    \theutterance \stepcounter{utterance}  
    & & & \multicolumn{2}{p{0.3\linewidth}}{
        \cellcolor[rgb]{0.9,0.9,0.9}{
            \makecell[{{p{\linewidth}}}]{
                \texttt{\tiny{[GM$|$GM]}}
                \texttt{yes} \\
            }
        }
    }
    & & \\ \\

    \theutterance \stepcounter{utterance}  
    & & & \multicolumn{2}{p{0.3\linewidth}}{
        \cellcolor[rgb]{0.9,0.9,0.9}{
            \makecell[{{p{\linewidth}}}]{
                \texttt{\tiny{[GM$|$GM]}}
                \texttt{game\_result = WIN} \\
            }
        }
    }
    & & \\ \\

\end{supertabular}
}

\end{document}
